\section{Related Work}
\label{sec:related_work}

\textbf{Static Model Merging.} Weight-space model merging unifies specialized networks into a single multi-task model without joint re-training. Early methods used simple arithmetic averaging \cite{wortsman2022model}, while task arithmetic \cite{ilharco2022editing} showed that fine-tuning trajectories act as ``task vectors'' that can be added or subtracted to modify behaviors. To resolve representational clash, recent works introduced TIES-merging \cite{yadav2023ties} (resolving sign conflicts) and Fisher-weighted averaging \cite{matena2022merging}. Other approaches address architectural alignment modulo permutation symmetries, such as Git Re-Basin \cite{ainsworth2022git, ainsworth2022linear} and permutation alignment \cite{pena2023permutation}, or representation renormalization like REPAIR \cite{jordan2022repair} and ZipIt! \cite{stoica2023zipit, stoica2024merging}. While parameter-efficient, these offline static techniques apply a single, global, sample-independent merging configuration, forcing a static compromise that underperforms under high task conflict.

\textbf{Dynamic Model Merging and Mixture-of-Experts (MoE).} To overcome static limits, dynamic model merging computes sample-specific merging coefficients on the fly \cite{gu2024zero, yadav2024resolving}, sharing key principles with Mixture-of-Experts (MoE) \cite{shazeer2017outrageously, fedus2022switch}. However, dynamic merging on the fly introduces major systems-level bottlenecks. Merging large parameters for input-dependent batches leads to severe memory-bandwidth overhead that can make dynamic merging slower than running independent models. Furthermore, under low routing temperatures, routing becomes discrete, rendering dynamic weight blending a systems tautology compared to direct sample routing.

\textbf{The Layer-Averaging Collapse Debate.} A critical design dimension of dynamic merging is the spatial resolution of routing. While layer-wise routing learns a specialized policy for each layer, a prominent theoretical study \cite{anonymous} proved a theorem of ``Layer-Averaging Collapse'' (rank-1 collapse), asserting that learned layer-wise routing coefficients are mathematically redundant and inevitably collapse to a collinear rank-1 subspace. 
However, their empirical validation was conducted almost exclusively within a simplified, 14-layer linear representation-space ``sandbox'' where weight-space interpolation is mathematically identical to taking a weighted linear combination of output logits. This sandbox completely hides the non-linear, high-dimensional weight-space dynamics of deeply hierarchical networks where layers extract highly distinct semantic abstractions.

\textbf{Representation Analysis in Deep Networks.} Landmark frameworks like SVCCA \cite{raghu2017svcca}, Projection-Weighted CCA \cite{morcos2018insights}, and CKA \cite{kornblith2019similarity} analyze representational drift and alignment across models \cite{conwell2021what, sucholutsky2023getting}. In this work, we adapt Singular Value Decomposition (SVD) and pairwise inter-layer cosine similarity to analyze learned dynamic routing coefficient matrices, establishing empirical boundaries where prior collinearity assumptions break down.
